\documentclass[11pt]{letter}
\usepackage[margin=1in]{geometry}
\usepackage{hyperref}

\signature{Alia Wu\\
Redline Rising \& Risk Efficacy\\
\texttt{wut08@nyu.edu}\\
ORCID: 0009-0005-4424-102X}

\address{Alia Wu\\Redline Rising\\wut08@nyu.edu}

\date{February 23, 2026}

\begin{document}

\begin{letter}{The Editors\\Nature Human Behaviour}

\opening{Dear Editors,}

I am submitting ``Cognitive architecture moderates emotion--cognition
relationships: cross-domain evidence from meta-analytic reanalysis'' for
consideration as an Article in \textit{Nature Human Behaviour}.

\textbf{Referral.} This manuscript was referred to
\textit{Nature Human Behaviour} by Mary Elizabeth at
\textit{Nature}.

\textbf{Summary.} The paper addresses a longstanding question:
why do meta-analyses of emotion--cognition relationships show
persistent heterogeneity that conventional moderators fail to
resolve? We propose that cognitive architecture---formalised by the
Dot--Linear--Network (DLN) framework---functions as a hidden
structural moderator. Across five moderator analyses spanning four
domains (emotion regulation, interoception, health psychology,
implicit cognition), DLN stage coding reduced between-study
heterogeneity ($\hat\tau^2$) by 70--100\% in four category-level
analyses and by 32.9\% at the study level in the largest dataset
($k = 184$, implicit cognition).
Exhaustive permutation tests confirmed this reduction is
non-arbitrary (combined $p < 0.001$). Critically, one analysis
revealed a dangerous-middle pattern---partial integration without
metacognitive monitoring produced worse outcomes than full
compartmentalisation---a direction-specific prediction unique to
the DLN framework. A boundary-condition dataset correctly predicted
null moderation.

\textbf{Why NHB.} The work (1) presents original empirical
results from meta-analytic reanalysis across four domains of
psychology (emotion regulation, interoception, health psychology,
implicit cognition), (2) demonstrates a cross-domain framework
with pre-specified replication rubrics, and (3) identifies a
computationally derived prediction (the dangerous-middle effect)
with direct implications for intervention design---a scope that
matches NHB's mandate for interdisciplinary behavioural science.

\textbf{Methodological strengths.} Pre-specified coding rubrics
for two additional meta-analytic datasets (Greenwald et al., 2009;
Zanini et al., 2025) are deposited in the public code repository,
enabling independent replication before and after publication. The
complete analysis pipeline, including permutation tests and REML
validation, is publicly available.

\textbf{Reproducibility.} All coding rubrics, extraction files,
analysis scripts, and permutation tests are deposited in a public
repository, enabling any researcher to independently verify every
result reported in the manuscript. A systematic coding sensitivity
analysis confirms that results survive all plausible alternative
stage assignments for debatable cases.

\textbf{Double-blind review.} This manuscript is submitted for
double-blind peer review. Author identifying information has been
removed from the manuscript file. Author name, affiliations, and
contact details are as follows: Alia Wu; Redline Rising (501(c)(3)
nonprofit) and Risk Efficacy; wut08@nyu.edu; ORCID
0009-0005-4424-102X.

\textbf{Anonymised citation.} For double-blind review, one
self-citation has been anonymised in the reference list and
manuscript text. The full citation is: Wu, A.\ (2026). Compression
Efficiency and Structural Learning as a Computational Model of DLN
Cognitive Stages. \textit{bioRxiv},
\url{https://doi.org/10.64898/2026.02.01.703168}.

\textbf{Related manuscript.} The companion paper presenting the DLN
computational framework (cited above) is currently under
consideration at \textit{PNAS}. The present manuscript is an
independent empirical test of that framework and does not overlap in
data, methods, or text.

\textbf{Competing interests.} A.W.\ is the President of Redline
Rising, a 501(c)(3) nonprofit, and founder of Risk Efficacy. The DLN
framework was developed as part of independent research and is not a
commercial product. The author declares no financial competing
interests.

\textbf{Author contributions.} Alia Wu: Conceptualization,
Methodology, Formal Analysis, Investigation, Data Curation,
Writing -- Original Draft, Writing -- Review \& Editing,
Visualization, Software.

\textbf{Acknowledgements.} The author thanks Yuanxi Jiang and Ang Li
for serving as independent blinded coders for the inter-rater
reliability validation, and the developers of the open-source tools
used in this analysis (numpy, scipy, matplotlib). The author also
thanks the original authors of the meta-analyses reanalysed here for
making their data accessible.

\textbf{Data and code availability.} All data, analysis scripts,
coding rubrics, and the complete permutation test are publicly
available at \url{https://github.com/aliawu08/dln-emotion-cognition}.
This URL has been omitted from the manuscript for double-blind review.

\textbf{Preprint disclosure.} A preprint of the present manuscript
is deposited at PsychArchives (10.23668/psycharchives.21641).

I confirm that this work has not been published elsewhere and is
not under consideration at any other journal. The related companion
paper noted above is under review at \textit{PNAS}.

\closing{Sincerely,}

\end{letter}
\end{document}
